% !TeX program = uplatex
\documentclass[a4paper,12pt]{jlreq}

% ===== Packages =====
\usepackage{amsmath,amssymb,mathtools,bm}
\usepackage{siunitx}
\sisetup{detect-all,per-mode=symbol}
\usepackage{physics}
\usepackage{mhchem}
\usepackage{booktabs}
\usepackage{graphicx}
\usepackage{hyperref}
\usepackage{ascmac}
\usepackage{xcolor}


\hypersetup{
  colorlinks=true,
  linkcolor=black,
  urlcolor=blue,
  citecolor=black
}
\usepackage{enumitem}
\usepackage{geometry}
\usepackage{fancybox}
\geometry{margin=25mm}

\begin{document}

\paragraph{問題 1.1.1 解答}

閉鎖系にある一定量の理想気体の体積変化に伴う仕事 $w$ は，
外圧を $p_{\mathrm{ex}}$ とすると
\[
  dw = -p_{\mathrm{ex}}\,dV
\]
で与えられる．

%%%%%%%%%%%%%%%%%%%%%%%%%%%%%%%%%%%%%%%%%%%%%%%%%%%%%%%%%%%%
\subparagraph{(1)}
初期状態と最終状態は
\[
  V_1 = 3.0\times 10^{-3}\ \mathrm{m^3},\quad
  p_1 = 1.0\times 10^{5}\ \mathrm{Pa},
\]
\[
  V_2 = 1.0\times 10^{-3}\ \mathrm{m^3},\quad
  p_2 = 3.0\times 10^{5}\ \mathrm{Pa}
\]
であり，$V_2<V_1$ なので圧縮過程である．
温度一定かつ理想気体なので $pV=nRT$ より
\[
  nRT = p_1 V_1 = (1.0\times10^{5})(3.0\times10^{-3})
  = 3.0\times10^{2}\ \mathrm{J}
\]
となる．

\begin{enumerate}
  \item[(i)] \textbf{外圧 $p_2$ による 1 段階圧縮}

  外圧は一定 $p_{\mathrm{ex}} = p_2$ であるから
  \[
    w = -\int_{V_1}^{V_2} p_{\mathrm{ex}}\,dV
      = -p_2(V_2-V_1)
      = -3.0\times10^{5}\,(1.0-3.0)\times10^{-3}.
  \]
  体積差 $V_2-V_1 = -2.0\times10^{-3}\ \mathrm{m^3}$ を代入すると
  \[
    w = -3.0\times10^{5}\times(-2.0\times10^{-3})
      = 6.0\times10^{2}\ \mathrm{J}.
  \]
  圧縮なので $w>0$ になっていることも確認できる．

  \item[(ii)] \textbf{準静的（可逆）等温圧縮}

  準静的過程では常に系の内部圧 $p$ が外圧に等しいので，
  \[
    pV = nRT = \text{一定}
  \]
  から
  \[
    p = \frac{nRT}{V}
  \]
  となる．したがって仕事は
  \[
    w = -\int_{V_1}^{V_2} p\,dV
      = -\int_{V_1}^{V_2} \frac{nRT}{V}\,dV
      = -nRT\int_{V_1}^{V_2} \frac{dV}{V}
      = -nRT \ln\frac{V_2}{V_1}.
  \]
  ここで $nRT = p_1V_1 = 3.0\times10^{2}\ \mathrm{J}$，
  \(
    V_2/V_1 = (1.0\times10^{-3})/(3.0\times10^{-3}) = 1/3
  \)
  なので
  \[
    w = -3.0\times10^{2}\ln\frac{1}{3}
      = 3.0\times10^{2}\ln 3
      \simeq 3.3\times10^{2}\ \mathrm{J}.
  \]
\end{enumerate}

%%%%%%%%%%%%%%%%%%%%%%%%%%%%%%%%%%%%%%%%%%%%%%%%%%%%%%%%%%%%
\subparagraph{(2)}
次に
\[
  V_1 = 2.1\ \mathrm{L},\quad p_1 = 2.0\ \mathrm{bar},
\]
\[
  V_2 = 3.0\ \mathrm{L},\quad p_2 = 1.4\ \mathrm{bar}
\]
であり，$V_2>V_1$ なので膨張過程である．
温度一定より
\[
  nRT = p_1V_1 = (2.0)(2.1) = 4.2\ \mathrm{L\cdot bar}.
\]
（計算の最後に $1\ \mathrm{L\cdot bar}=100\ \mathrm{J}$ を用いる．）

\begin{enumerate}
  \item[(i)] \textbf{外圧 $p_2$ による 1 段階膨張}

  外圧一定 $p_{\mathrm{ex}} = p_2$ なので
  \[
    w = -p_{\mathrm{ex}}(V_2-V_1)
      = -1.4\,(3.0-2.1)\ \mathrm{L\cdot bar}
      = -1.26\ \mathrm{L\cdot bar}.
  \]
  これをジュールに直すと
  \[
    w = -1.26\times 100\ \mathrm{J}
      \simeq -1.3\times10^{2}\ \mathrm{J}.
  \]

  \item[(ii)] \textbf{準静的（可逆）等温膨張}

  (1) と同様に
  \[
    w = -\int_{V_1}^{V_2} p\,dV
      = -nRT\ln\frac{V_2}{V_1}.
  \]
  ここで $nRT = 4.2\ \mathrm{L\cdot bar}$，
  \(
    V_2/V_1 = 3.0/2.1 = 10/7
  \)
  だから
  \[
    w = -4.2\ln\!\left(\frac{10}{7}\right)\ \mathrm{L\cdot bar}
      \simeq -1.41\ \mathrm{L\cdot bar}.
  \]
  ジュールに変換すると
  \[
    w \simeq -1.41\times100\ \mathrm{J}
      \simeq -1.4\times10^{2}\ \mathrm{J}.
  \]
  膨張過程なので $w<0$ となっており，符号規約とも一致している．
\end{enumerate}






\newpage


\paragraph{1.1.2 の解答（メタンの燃焼）}

反応：
\[
\mathrm{CH_4(g) + 2O_2(g) \rightarrow CO_2(g) + 2H_2O(l)}
\]

標準生成エンタルピー：
\[
\Delta_{\mathrm f}H^\circ(\mathrm{CH_4}) = -74.81\ \mathrm{kJ\,mol^{-1}},
\quad
\Delta_{\mathrm f}H^\circ(\mathrm{CO_2}) = -393.51\ \mathrm{kJ\,mol^{-1}},
\]
\[
\Delta_{\mathrm f}H^\circ(\mathrm{H_2O(l)}) = -285.83\ \mathrm{kJ\,mol^{-1}},
\quad
\Delta_{\mathrm f}H^\circ(\mathrm{O_2(g)}) = 0.
\]

標準反応エンタルピーの一般式
\[
\Delta_{\mathrm r}H^\circ
 = \sum_{\text{生成系}} \nu_i \Delta_{\mathrm f}H_i^\circ
 - \sum_{\text{反応原系}} \nu_i \Delta_{\mathrm f}H_i^\circ
\]
に従うと，

生成系：
\[
\sum_{\text{生成系}} \nu_i \Delta_{\mathrm f}H_i^\circ
= 1\cdot\Delta_{\mathrm f}H^\circ(\mathrm{CO_2})
  + 2\cdot\Delta_{\mathrm f}H^\circ(\mathrm{H_2O(l)})
\]
\[
= (-393.51) + 2(-285.83)
= -393.51 - 571.66
= -965.17\ \mathrm{kJ\,mol^{-1}}.
\]

反応原系：
\[
\sum_{\text{反応原系}} \nu_i \Delta_{\mathrm f}H_i^\circ
= 1\cdot\Delta_{\mathrm f}H^\circ(\mathrm{CH_4})
  + 2\cdot\Delta_{\mathrm f}H^\circ(\mathrm{O_2})
= -74.81 + 2\cdot 0
= -74.81\ \mathrm{kJ\,mol^{-1}}.
\]

したがって
\[
\Delta_{\mathrm r}H^\circ
= -965.17 - (-74.81)
= -890.36\ \mathrm{kJ\,mol^{-1}}.
\]

-------------------------------------------------------------------------------

\paragraph{1.1.3 の解答}

反応：
\[
\mathrm{O_2(g) + 4HCl(g) \rightarrow 2Cl_2(g) + 2H_2O(g)}
\]

与えられた標準生成エンタルピー：
\[
\Delta_{\mathrm f}H^\circ(\mathrm{HCl(g)}) = -92.31\ \mathrm{kJ\,mol^{-1}},
\quad
\Delta_{\mathrm f}H^\circ(\mathrm{H_2O(g)}) = -241.82\ \mathrm{kJ\,mol^{-1}}.
\]
元素の標準状態については
\[
\Delta_{\mathrm f}H^\circ(\mathrm{O_2(g)}) = 0,\quad
\Delta_{\mathrm f}H^\circ(\mathrm{Cl_2(g)}) = 0.
\]

生成系：
\[
\sum_{\text{生成系}} \nu_i \Delta_{\mathrm f}H_i^\circ
= 2\cdot\Delta_{\mathrm f}H^\circ(\mathrm{Cl_2(g)})
  + 2\cdot\Delta_{\mathrm f}H^\circ(\mathrm{H_2O(g)})
\]
\[
= 2\cdot 0 + 2(-241.82)
= -483.64\ \mathrm{kJ\,mol^{-1}}.
\]

反応原系：
\[
\sum_{\text{反応原系}} \nu_i \Delta_{\mathrm f}H_i^\circ
= 1\cdot\Delta_{\mathrm f}H^\circ(\mathrm{O_2(g)})
  + 4\cdot\Delta_{\mathrm f}H^\circ(\mathrm{HCl(g)})
\]
\[
= 0 + 4(-92.31)
= -369.24\ \mathrm{kJ\,mol^{-1}}.
\]

したがって
\[
\Delta_{\mathrm r}H^\circ
= -483.64 - (-369.24)
= -114.40\ \mathrm{kJ\,mol^{-1}}.
\]

-------------------------------------------------------------------------------

\paragraph{1.1.4 の解答}

反応：
\[
\mathrm{Fe_2O_3(c) + 3CO(g) \rightarrow 2Fe(c) + 3CO_2(g)}
\]

標準生成エンタルピー（表より）：
\[
\Delta_{\mathrm f}H^\circ(\mathrm{Fe_2O_3(c)}) = -824.2\ \mathrm{kJ\,mol^{-1}},
\quad
\Delta_{\mathrm f}H^\circ(\mathrm{CO(g)}) = -110.5\ \mathrm{kJ\,mol^{-1}},
\]
\[
\Delta_{\mathrm f}H^\circ(\mathrm{CO_2(g)}) = -393.5\ \mathrm{kJ\,mol^{-1}}.
\]
元素の固体鉄は標準状態なので
\[
\Delta_{\mathrm f}H^\circ(\mathrm{Fe(c)}) = 0.
\]

生成系：
\[
\sum_{\text{生成系}} \nu_i \Delta_{\mathrm f}H_i^\circ
= 2\cdot\Delta_{\mathrm f}H^\circ(\mathrm{Fe(c)})
  + 3\cdot\Delta_{\mathrm f}H^\circ(\mathrm{CO_2(g)})
\]
\[
= 2\cdot 0 + 3(-393.5)
= -1180.5\ \mathrm{kJ\,mol^{-1}}.
\]

反応原系：
\[
\sum_{\text{反応原系}} \nu_i \Delta_{\mathrm f}H_i^\circ
= 1\cdot\Delta_{\mathrm f}H^\circ(\mathrm{Fe_2O_3(c)})
  + 3\cdot\Delta_{\mathrm f}H^\circ(\mathrm{CO(g)})
\]
\[
= (-824.2) + 3(-110.5)
= -824.2 - 331.5
= -1155.7\ \mathrm{kJ\,mol^{-1}}.
\]

したがって
\[
\Delta_{\mathrm r}H^\circ
= -1180.5 - (-1155.7)
= -24.8\ \mathrm{kJ\,mol^{-1}}.
\]

-------------------------------------------------------------------------------

\noindent
よって，求める標準反応エンタルピーは
\[
\boxed{\Delta_{\mathrm r}H^\circ(1.1.2) = -8.90\times10^{2}\ \mathrm{kJ\,mol^{-1}}}
\]
\[
\boxed{\Delta_{\mathrm r}H^\circ(1.1.3) = -1.14\times10^{2}\ \mathrm{kJ\,mol^{-1}}}
\]
\[
\boxed{\Delta_{\mathrm r}H^\circ(1.1.4) = -2.48\times10^{1}\ \mathrm{kJ\,mol^{-1}}}
\]
である。






\newpage




\paragraph{問題 1.3.1 解答}

講義ノートより，標準反応 Gibbs エネルギーは状態関数であり，
標準生成 Gibbs エネルギー $\Delta_{\mathrm f}G_i^\circ$ を用いて
\[
\Delta_{\mathrm r}G^\circ
  = \sum_{\text{生成系}} \nu_i \Delta_{\mathrm f}G_i^\circ
  - \sum_{\text{反応原系}} \nu_i \Delta_{\mathrm f}G_i^\circ
\]
と表される．（$\nu_i$ は化学量論係数）:contentReference[oaicite:0]{index=0}  

$\Delta_{\mathrm r}G^\circ<0$ なら標準状態で自発的に進行し，
$\Delta_{\mathrm r}G^\circ>0$ なら自発的に進行しない。

表に与えられた標準生成 Gibbs エネルギー：
\[
\begin{array}{lcl}
\Delta_{\mathrm f}G^\circ(\mathrm{CO(g)}) &=& -137.17\ \mathrm{kJ\,mol^{-1}}\\
\Delta_{\mathrm f}G^\circ(\mathrm{CO_2(g)}) &=& -394.36\ \mathrm{kJ\,mol^{-1}}\\
\Delta_{\mathrm f}G^\circ(\mathrm{NH_3(g)}) &=& -16.45\ \mathrm{kJ\,mol^{-1}}\\
\Delta_{\mathrm f}G^\circ(\mathrm{NO(g)})  &=& +86.55\ \mathrm{kJ\,mol^{-1}}\\
\Delta_{\mathrm f}G^\circ(\mathrm{MgO(s)}) &=& -569.45\ \mathrm{kJ\,mol^{-1}}\\[2mm]
\Delta_{\mathrm f}G^\circ(\mathrm{PbO(s)}) &=& -188.93\ \mathrm{kJ\,mol^{-1}}\\
\Delta_{\mathrm f}G^\circ(\mathrm{CaF_2(s)}) &=& -1167.3\ \mathrm{kJ\,mol^{-1}}\\
\Delta_{\mathrm f}G^\circ(\mathrm{Ca^{2+}(aq)}) &=& -553.58\ \mathrm{kJ\,mol^{-1}}\\
\Delta_{\mathrm f}G^\circ(\mathrm{F^-(aq)}) &=& -278.79\ \mathrm{kJ\,mol^{-1}}
\end{array}
\]
元素の標準状態については
\[
\Delta_{\mathrm f}G^\circ(\mathrm{Pb(s)}) =
\Delta_{\mathrm f}G^\circ(\mathrm{Mg(s)}) =
\Delta_{\mathrm f}G^\circ(\mathrm{N_2(g)}) =
\Delta_{\mathrm f}G^\circ(\mathrm{O_2(g)}) =
\Delta_{\mathrm f}G^\circ(\mathrm{H_2(g)}) = 0
\]
である。

%----------------------------------------------------------
\subparagraph{(1)\ $\mathrm{PbO(s)+CO(g)\rightleftharpoons Pb(s)+CO_2(g)}$}

生成系：
\[
\sum_{\text{生成系}} \nu_i \Delta_{\mathrm f}G_i^\circ
= \Delta_{\mathrm f}G^\circ(\mathrm{Pb(s)}) + \Delta_{\mathrm f}G^\circ(\mathrm{CO_2(g)})
= 0 + (-394.36)
= -394.36
\]

反応原系：
\[
\sum_{\text{反応原系}} \nu_i \Delta_{\mathrm f}G_i^\circ
= \Delta_{\mathrm f}G^\circ(\mathrm{PbO(s)}) + \Delta_{\mathrm f}G^\circ(\mathrm{CO(g)})
= -188.93 + (-137.17)
= -326.10
\]

したがって
\[
\Delta_{\mathrm r}G^\circ
= -394.36 - (-326.10)
= -68.26\ \mathrm{kJ\,mol^{-1}}.
\]

$\Delta_{\mathrm r}G^\circ<0$ より，\textbf{標準状態で自発的に進行する}。

%----------------------------------------------------------
\subparagraph{(2)\ $\mathrm{MgO(s)+CO(g)\rightleftharpoons Mg(s)+CO_2(g)}$}

生成系：
\[
\sum_{\text{生成系}} \nu_i \Delta_{\mathrm f}G_i^\circ
= \Delta_{\mathrm f}G^\circ(\mathrm{Mg(s)}) + \Delta_{\mathrm f}G^\circ(\mathrm{CO_2(g)})
= 0 + (-394.36)
= -394.36
\]

反応原系：
\[
\sum_{\text{反応原系}} \nu_i \Delta_{\mathrm f}G_i^\circ
= \Delta_{\mathrm f}G^\circ(\mathrm{MgO(s)}) + \Delta_{\mathrm f}G^\circ(\mathrm{CO(g)})
= -569.45 + (-137.17)
= -706.62
\]

したがって
\[
\Delta_{\mathrm r}G^\circ
= -394.36 - (-706.62)
= +312.26\ \mathrm{kJ\,mol^{-1}}.
\]

$\Delta_{\mathrm r}G^\circ>0$ なので，\textbf{標準状態では自発的に進行しない}
（逆反応の方が自発的）。

%----------------------------------------------------------
\subparagraph{(3)\ $\mathrm{N_2(g)+O_2(g)\rightleftharpoons 2NO(g)}$}

生成系：
\[
\sum_{\text{生成系}} \nu_i \Delta_{\mathrm f}G_i^\circ
= 2\,\Delta_{\mathrm f}G^\circ(\mathrm{NO(g)})
= 2\times 86.55
= 173.10
\]

反応原系：
\[
\sum_{\text{反応原系}} \nu_i \Delta_{\mathrm f}G_i^\circ
= \Delta_{\mathrm f}G^\circ(\mathrm{N_2(g)})
 + \Delta_{\mathrm f}G^\circ(\mathrm{O_2(g)})
= 0+0=0
\]

したがって
\[
\Delta_{\mathrm r}G^\circ = 173.10\ \mathrm{kJ\,mol^{-1}}.
\]

$\Delta_{\mathrm r}G^\circ>0$ より，
\textbf{標準状態では自発的に進行しない}。

%----------------------------------------------------------
\subparagraph{(4)\ $\mathrm{CaF_2(s)\rightleftharpoons Ca^{2+}(aq)+2F^-(aq)}$}

生成系：
\[
\sum_{\text{生成系}} \nu_i \Delta_{\mathrm f}G_i^\circ
= \Delta_{\mathrm f}G^\circ(\mathrm{Ca^{2+}(aq)})
  + 2\,\Delta_{\mathrm f}G^\circ(\mathrm{F^-(aq)})
= -553.58 + 2(-278.79)
= -1111.16
\]

反応原系：
\[
\sum_{\text{反応原系}} \nu_i \Delta_{\mathrm f}G_i^\circ
= \Delta_{\mathrm f}G^\circ(\mathrm{CaF_2(s)})
= -1167.3
\]

したがって
\[
\Delta_{\mathrm r}G^\circ
= -1111.16 - (-1167.3)
= +56.14\ \mathrm{kJ\,mol^{-1}}.
\]

よって $\Delta_{\mathrm r}G^\circ>0$ であり，
\textbf{標準状態では溶解は自発的に進行しない}
（固体 $\mathrm{CaF_2}$ のままの方が安定）。

%----------------------------------------------------------
\subparagraph{(5)\ $\mathrm{N_2(g)+3H_2(g)\rightleftharpoons 2NH_3(g)}$}

生成系：
\[
\sum_{\text{生成系}} \nu_i \Delta_{\mathrm f}G_i^\circ
= 2\,\Delta_{\mathrm f}G^\circ(\mathrm{NH_3(g)})
= 2\times(-16.45)
= -32.90
\]

反応原系：
\[
\sum_{\text{反応原系}} \nu_i \Delta_{\mathrm f}G_i^\circ
= \Delta_{\mathrm f}G^\circ(\mathrm{N_2(g)})
  + 3\,\Delta_{\mathrm f}G^\circ(\mathrm{H_2(g)})
= 0 + 3\times 0 = 0
\]

したがって
\[
\Delta_{\mathrm r}G^\circ
= -32.90\ \mathrm{kJ\,mol^{-1}}.
\]

$\Delta_{\mathrm r}G^\circ<0$ なので，
\textbf{標準状態でアンモニア生成反応は自発的に進行する}。

%----------------------------------------------------------
\noindent
以上をまとめると，
\[
\begin{array}{c|c|c}
\text{反応} & \Delta_{\mathrm r}G^\circ\ /\ \mathrm{kJ\,mol^{-1}} & \text{標準状態での自発性} \\
\hline
(1) & -68.26  & \text{自発的} \\
(2) & +312.26 & \text{自発的でない（逆向きが自発的）} \\
(3) & +173.10 & \text{自発的でない} \\
(4) & +56.14  & \text{自発的でない} \\
(5) & -32.90  & \text{自発的}
\end{array}
\]
となる。





\newpage



\paragraph{問題 1.3.2 解答}

KOH 型燃料電池
\[
\mathrm{H_2(Pt)\,|\,KOH(aq)\,|\,O_2(Pt)}
\]
では，各電極で
\[
\begin{cases}
(-)\ \mathrm{H_2 + 2OH^- \rightarrow 2H_2O(l) + 2e^-} \\
(+)\ \dfrac{1}{2}\mathrm{O_2 + H_2O(l) + 2e^- \rightarrow 2OH^-}
\end{cases}
\]
の電極反応が起こる。ギブズエネルギーと最大外部仕事の関係は
\[
-\Delta_{\mathrm r}G^\circ = -w_{e,\max}
\]
であり，さらに
\[
\Delta_{\mathrm r}G^\circ = -\nu F E^\circ
\]
が成り立つ。ここで $\nu$ は電子の量論数，$F$ は Faraday 定数である:contentReference[oaicite:0]{index=0}。

\subparagraph{(1) 全体の酸化還元反応}

2e$^-$，OH$^-$，H$_2$O を消去して全体の反応式を求める：
\[
\mathrm{H_2 + \frac{1}{2}O_2 \rightarrow H_2O(l)}.
\]

\subparagraph{(2) 4.00 mol の水素をすべて消費したときの最大仕事}

上で得た全体反応は，
標準状態における水の生成反応
\[
\mathrm{H_2(g) + \frac{1}{2}O_2(g) \rightarrow H_2O(l)}
\]
そのものであるから，
\[
\Delta_{\mathrm r}G^\circ
  = \Delta_{\mathrm f}G^\circ(\mathrm{H_2O(l)})
  - \Delta_{\mathrm f}G^\circ(\mathrm{H_2})
  - \frac{1}{2}\Delta_{\mathrm f}G^\circ(\mathrm{O_2})
  = -237.1\ \mathrm{kJ\,mol^{-1}}.
\]

4.00 mol の H$_2$ が反応するので，全反応のギブズエネルギー変化は
\[
\Delta G = 4.00\times(-237.1)\ \mathrm{kJ}
          = -948.4\ \mathrm{kJ}.
\]
ギブズエネルギーの減少量の絶対値が外部に取り出せる最大非膨張仕事であるから，
\[
-w_{e,\max} = \Delta G \quad\Rightarrow\quad
w_{e,\max} = -\Delta G = 9.48\times10^{2}\ \mathrm{kJ}.
\]

\subparagraph{(3) この電池の標準起電力 $E^\circ$}

全体反応 1 回あたりに移動する電子数は，電極反応から
\[
\nu = 2
\]
である。よって
\[
E^\circ
 = -\frac{\Delta_{\mathrm r}G^\circ}{\nu F}
 = -\frac{-237.1\times10^{3}\ \mathrm{J\,mol^{-1}}}
        {2\times 9.6485\times10^{4}\ \mathrm{C\,mol^{-1}}}.
\]
計算すると
\[
E^\circ \simeq 1.23\ \mathrm{V}.
\]

\subparagraph{(4) リン酸型燃料電池の標準起電力}

リン酸型燃料電池
\[
\mathrm{H_2(Pt)\,|\,H_3PO_4(aq)\,|\,O_2(Pt)}
\]
では，各電極反応は
\[
\begin{cases}
(-)\ \mathrm{H_2 \rightarrow 2H^+ + 2e^-} \\
(+)\ \dfrac{1}{2}\mathrm{O_2 + 2H^+ + 2e^- \rightarrow H_2O(l)}
\end{cases}
\]
である。両辺を加えると
\[
\mathrm{H_2 + \frac{1}{2}O_2 \rightarrow H_2O(l)}
\]
となり，全体反応は KOH 型燃料電池と同じである。
したがって $\Delta_{\mathrm r}G^\circ$ も同じ値をとり，
電子数 $\nu = 2$ であることも変わらないので，
\[
E^\circ = 1.23\ \mathrm{V}
\]
と求まる。

\paragraph{問題 1.3.3 解答}

鉛蓄電池
\[
\mathrm{Pb\,|\,H_2SO_4(aq)\,|\,PbO_2}
\]
の標準起電力は $E^\circ = 2.05\ \mathrm{V}$ であり，起こっている全体の酸化還元反応は
\[
\mathrm{PbO_2 + Pb + 2H_2SO_4 \rightarrow 2PbSO_4 + 2H_2O}
\]
である。

半反応は例えば
\[
\mathrm{Pb + SO_4^{2-} \rightarrow PbSO_4 + 2e^-}
\]
\[
\mathrm{PbO_2 + SO_4^{2-} + 4H^+ + 2e^- \rightarrow PbSO_4 + 2H_2O}
\]
のように書けるので，全体反応 1 回あたりの電子の量論数は
\[
\nu = 2
\]
である。

したがって標準反応ギブズエネルギーは
\[
\Delta_{\mathrm r}G^\circ
= -\nu F E^\circ
= -2 \times 9.6485\times10^{4}\ \mathrm{C\,mol^{-1}}
     \times 2.05\ \mathrm{V}.
\]
\[
\Delta_{\mathrm r}G^\circ
\simeq -3.96\times10^{5}\ \mathrm{J\,mol^{-1}}
= -3.96\times10^{2}\ \mathrm{kJ\,mol^{-1}}.
\]

よって，
\[
\boxed{\Delta_{\mathrm r}G^\circ \approx -4.0\times10^{2}\ \mathrm{kJ\,mol^{-1}}}
\]
となり，非常に自発的な反応であることがわかる。






%%%%%%%%%%%%%%%%%%%%%%%%%%%%%



\newpage

%%%%%%%%%%%%%%%%%%%%%%%%%%%%%%%%%%%%%%%%%%%%%%%%
% 1.4.1
%%%%%%%%%%%%%%%%%%%%%%%%%%%%%%%%%%%%%%%%%%%%%%%%
\subsection*{1.4.1 圧力変化に伴うギブズエネルギー変化}

温度一定・物質量一定では
\[
 \dd G = V\,\dd P
\]
であるから，モルギブズエネルギーの変化は
\[
 \Delta G = \int_{P_1}^{P_2} \bar V\,\dd P
\]
で与えられる．

\paragraph{(1) 氷の圧縮}

氷の密度を $\rho = 917\ \mathrm{kg\,m^{-3}}$，
水のモル質量を $M=18.0\times 10^{-3}\ \mathrm{kg\,mol^{-1}}$ とすると，
氷1 mol のモル体積は
\[
 \bar V = \frac{M}{\rho}
      = \frac{18.0\times10^{-3}}{917}
      \simeq 1.96\times10^{-5}\ \mathrm{m^3\,mol^{-1}}
\]
でほぼ一定とみなせる．  
圧力を $P_1=1\ \mathrm{bar}$ から $P_2=2\ \mathrm{bar}$ まで変化させるので
\[
 \Delta P = P_2-P_1 = 1\ \mathrm{bar}
          = 1.0\times10^{5}\ \mathrm{Pa}
\]
より
\[
 \Delta G
 \simeq \bar V\,\Delta P
 = (1.96\times10^{-5})(1.0\times10^{5})
 \simeq 2.0\ \mathrm{J\,mol^{-1}}.
\]
したがって，氷1 mol のギブズエネルギーはおよそ
\[
 \boxed{\Delta G \simeq 2\ \mathrm{J\,mol^{-1}}}
\]
だけ増加する．

\paragraph{(2) 水蒸気の圧縮（理想気体）}

理想気体ではモルギブズエネルギーは
\[
 G_m(T,P) = G_m^\circ(T) + RT\ln\frac{P}{P^\circ}
\]
で与えられるので，温度一定で圧力を
$P_1=1\ \mathrm{bar}$ から $P_2=2\ \mathrm{bar}$ に変化させたとき
\[
 \Delta G = G_m(T,P_2)-G_m(T,P_1)
          = RT\ln\frac{P_2}{P_1}
          = RT\ln2
\]
となる．$T=25^\circ\mathrm{C}=298.15\ \mathrm{K}$，
$R=8.314\ \mathrm{J\,mol^{-1}K^{-1}}$ を代入すると
\[
 \Delta G
 = (8.314)(298.15)\ln2
 \simeq 1.7\times10^{3}\ \mathrm{J\,mol^{-1}}
\]
したがって
\[
 \boxed{\Delta G \simeq 1.7\times10^{3}\ \mathrm{J\,mol^{-1}}}
\]
だけ増加する．

%%%%%%%%%%%%%%%%%%%%%%%%%%%%%%%%%%%%%%%%%%%%%%%%
% 1.4.2
%%%%%%%%%%%%%%%%%%%%%%%%%%%%%%%%%%%%%%%%%%%%%%%%
\subsection*{1.4.2 量論数}

量論数 $\nu_i$ は，反応の進行度 $\xi$ に対して
\[
 \dd n_i = \nu_i\,\dd\xi
\]
を満たすように定義され，反応物では負，生成物では正となる．

\paragraph{(1)}
\[
 \frac{1}{2}\mathrm{N_2} + \frac{1}{2}\mathrm{O_2} \rightarrow \mathrm{NO}
\]
より
\[
 \boxed{
 \nu_{\mathrm{N_2}} = -\frac{1}{2},\quad
 \nu_{\mathrm{O_2}} = -\frac{1}{2},\quad
 \nu_{\mathrm{NO}} = +1
 }
\]

\paragraph{(2)}
\[
 \frac{1}{2}\mathrm{N_2} + \frac{3}{2}\mathrm{H_2} \rightarrow \mathrm{NH_3}
\]
より
\[
 \boxed{
 \nu_{\mathrm{N_2}} = -\frac{1}{2},\quad
 \nu_{\mathrm{H_2}} = -\frac{3}{2},\quad
 \nu_{\mathrm{NH_3}} = +1
 }
\]

%%%%%%%%%%%%%%%%%%%%%%%%%%%%%%%%%%%%%%%%%%%%%%%%
% 1.4.3
%%%%%%%%%%%%%%%%%%%%%%%%%%%%%%%%%%%%%%%%%%%%%%%%
\subsection*{1.4.3 平衡状態とギブズエネルギー}

定温・定圧のもとで，閉じた系のギブズエネルギー $G$ は
自発的な変化に対して減少し，平衡で極小値をとる．  
反応の進行度 $\xi$ を変数とみなすと，
\[
 \left(\frac{\partial G}{\partial \xi}\right)_{T,P} = \Delta_r G
\]
であり，平衡では
\[
 \boxed{
 \Delta_r G = 0
 \quad\Leftrightarrow\quad
 \left(\dfrac{\partial G}{\partial \xi}\right)_{T,P} = 0
 }
\]
となる．

したがって，図に示された $G$–$\xi$ 曲線においては，
\begin{itemize}
  \item 反応 A：$G$ が最小となる $\xi$ の値（谷の底の点）
  \item 反応 B：同様に，曲線 B の谷の底に対応する $\xi$ の値
\end{itemize}
がそれぞれの平衡組成を与える．  
（どちらも $G$–$\xi$ 曲線の極小点が平衡状態である。）


\newpage

%%%%%%%%%%%%%%%%%%%%%%%%%%%%%%%%%%%%%%%%%%%%%%%%
% 1.4.4
%%%%%%%%%%%%%%%%%%%%%%%%%%%%%%%%%%%%%%%%%%%%%%%%
\subsection*{1.4.4 反応ギブズエネルギーと自発的な方向}

反応：
\[
 \frac12\mathrm{N_2(g)} + \frac32\mathrm{H_2(g)} \rightleftharpoons \mathrm{NH_3(g)}
\]
の標準生成ギブズエネルギーは
\[
 \Delta_{\mathrm f}G^\circ(\mathrm{NH_3}) = -16.45\ \mathrm{kJ\,mol^{-1}}
\]
であり，\(\mathrm{N_2},\mathrm{H_2}\) の標準生成ギブズエネルギーは 0 であるから
\[
 \Delta_r G^\circ
 = \nu_{\mathrm{NH_3}}\Delta_{\mathrm f}G^\circ(\mathrm{NH_3})
   + \nu_{\mathrm{N_2}}\Delta_{\mathrm f}G^\circ(\mathrm{N_2})
   + \nu_{\mathrm{H_2}}\Delta_{\mathrm f}G^\circ(\mathrm{H_2})
 = -16.45\ \mathrm{kJ\,mol^{-1}}
\]
である．  

反応ギブズエネルギーは
\[
 \Delta_r G = \Delta_r G^\circ + RT\ln Q
\]
で与えられる．ここで反応比 \(Q\) は
\[
 Q = \prod_i\left(\frac{p_i}{p^\circ}\right)^{\nu_i}
\]
であり，この反応では
\[
 Q = \frac{(p_{\mathrm{NH_3}}/p^\circ)^1}
           {(p_{\mathrm{N_2}}/p^\circ)^{1/2}(p_{\mathrm{H_2}}/p^\circ)^{3/2}}
\]
となる．標準圧力を \(p^\circ = 1.0\times10^5\ \mathrm{Pa}\)（1 bar）とし，
与えられた分圧を
\[
 \tilde p_i = \frac{p_i}{p^\circ}
\]
とおくと，\(\tilde p\) は無次元量であり，計算が簡単になる．

温度は 298.15 K，気体定数は
\[
 R = 8.314\ \mathrm{J\,mol^{-1}K^{-1}}
\]
として数値計算する．

\paragraph{(1) 各分圧が \(0.30, 0.30, 0.40\times10^5\ \mathrm{Pa}\) のとき}

このとき
\[
 \tilde p_{\mathrm{N_2}} = 0.30,\quad
 \tilde p_{\mathrm{H_2}} = 0.30,\quad
 \tilde p_{\mathrm{NH_3}} = 0.40
\]
であるから
\[
 Q = \frac{0.40}{(0.30)^{1/2}(0.30)^{3/2}}
   = \frac{0.40}{(0.30)^2}
   = \frac{0.40}{0.09}
   \simeq 4.44.
\]
したがって
\[
 \Delta_r G
 = \Delta_r G^\circ + RT\ln Q
 = -16.45\times10^3
   + (8.314)(298.15)\ln(4.44)\ \mathrm{J\,mol^{-1}}.
\]
\(\ln(4.44)\simeq 1.49\) を用いると
\[
 RT\ln Q
 \simeq (8.314)(298.15)(1.49)
 \simeq 3.70\times10^3\ \mathrm{J\,mol^{-1}}
\]
より
\[
 \Delta_r G
 \simeq -1.645\times10^4 + 3.70\times10^3
 = -1.28\times10^4\ \mathrm{J\,mol^{-1}}
 = -12.8\ \mathrm{kJ\,mol^{-1}}.
\]
\(\Delta_r G<0\) であるから，この条件では
\[
 \boxed{\text{正反応（NH\(_3\) 生成向き）が自発的である}}
\]
と言える．

\paragraph{(2) 各分圧が \(0.01, 0.03, 0.96\times10^5\ \mathrm{Pa}\) のとき}

同様に
\[
 \tilde p_{\mathrm{N_2}} = 0.01,\quad
 \tilde p_{\mathrm{H_2}} = 0.03,\quad
 \tilde p_{\mathrm{NH_3}} = 0.96
\]
より
\[
 Q
 = \frac{0.96}{(0.01)^{1/2}(0.03)^{3/2}}
 \simeq 1.85\times10^3.
\]
したがって
\[
 \Delta_r G
 = -16.45\times10^3
   + (8.314)(298.15)\ln(1.85\times10^3)\ \mathrm{J\,mol^{-1}}.
\]
\(\ln(1.85\times10^3)\simeq 7.52\) より
\[
 RT\ln Q
 \simeq (8.314)(298.15)(7.52)
 \simeq 1.88\times10^4\ \mathrm{J\,mol^{-1}}
\]
なので
\[
 \Delta_r G
 \simeq -1.645\times10^4 + 1.88\times10^4
 = 2.19\times10^3\ \mathrm{J\,mol^{-1}}
 = 2.19\ \mathrm{kJ\,mol^{-1}}.
\]
\(\Delta_r G>0\) であるから，
\[
 \boxed{\text{この条件では正反応は自発的ではなく，逆反応（NH\(_3\) 分解向き）が自発的である}}
\]
ことが分かる．


%%%%%%%%%%%%%%%%%%%%%%%%%%%%%%%%%%%%%%%%%%%%%%%%
% 1.4.5
%%%%%%%%%%%%%%%%%%%%%%%%%%%%%%%%%%%%%%%%%%%%%%%%
\subsection*{1.4.5 N\texorpdfstring{$_2$}{2}O\texorpdfstring{$_4$}{4} の分解平衡}

反応：
\[
 \mathrm{N_2O_4(g)} \rightleftharpoons 2\mathrm{NO_2(g)}
\]

\paragraph{(1) 平衡状態における各分圧}

初めに \(\mathrm{N_2O_4}\) を 1 mol 封入したとし，分解の反応進行度を \(\xi\) とする．
\(\alpha\) を「分解した \(\mathrm{N_2O_4}\) の割合」とするとき，
問題文より \(\alpha = 0.1846 = 18.46~\mathrm{mol\%}\) であるから，
\[
 \xi = \alpha \times 1\ \mathrm{mol} = 0.1846\ \mathrm{mol}.
\]
平衡時の各成分の物質量は
\[
 n_{\mathrm{N_2O_4}} = 1 - \xi = 0.8154\ \mathrm{mol},
\qquad
 n_{\mathrm{NO_2}} = 2\xi = 0.3692\ \mathrm{mol}.
\]
全物質量は
\[
 n_{\mathrm{tot}} = n_{\mathrm{N_2O_4}} + n_{\mathrm{NO_2}}
 = 1.1846\ \mathrm{mol}
\]
である．したがってモル分率は
\[
 y_{\mathrm{N_2O_4}} = \frac{0.8154}{1.1846} \simeq 0.688,
\qquad
 y_{\mathrm{NO_2}} = \frac{0.3692}{1.1846} \simeq 0.312.
\]
全圧が \(P_{\mathrm{tot}} = 1.00\ \mathrm{bar}\)（= \(1.0\times10^5\ \mathrm{Pa}\)）なので，
平衡分圧は
\[
 \boxed{
 p_{\mathrm{N_2O_4}} \simeq 0.688\ \mathrm{bar},\qquad
 p_{\mathrm{NO_2}} \simeq 0.312\ \mathrm{bar}
 }.
\]

\paragraph{(2) 熱力学平衡定数と \(\Delta_r G^\circ\)}

平衡定数 \(K\) は（理想気体とし，標準圧力 \(p^\circ=1\ \mathrm{bar}\) とする）
\[
 K = \prod_i\left(\frac{p_i}{p^\circ}\right)^{\nu_i}
\]
より，
\[
 K = \frac{\left(p_{\mathrm{NO_2}}/p^\circ\right)^2}
          {p_{\mathrm{N_2O_4}}/p^\circ}
   = \frac{p_{\mathrm{NO_2}}^2}{p_{\mathrm{N_2O_4}}}
\]
で与えられる．したがって
\[
 K = \frac{(0.312)^2}{0.688}
   \simeq 0.141.
\]
標準反応ギブズエネルギーは
\[
 \Delta_r G^\circ = -RT\ln K
\]
で与えられるので，
\[
 \Delta_r G^\circ
 = -(8.314)(298.15)\ln(0.141)\ \mathrm{J\,mol^{-1}}
 \simeq 4.85\times10^3\ \mathrm{J\,mol^{-1}}
 = 4.85\ \mathrm{kJ\,mol^{-1}}.
\]
よって
\[
 \boxed{K \simeq 0.14,\qquad
        \Delta_r G^\circ \simeq 4.85\ \mathrm{kJ\,mol^{-1}}}
\]
である．


%%%%%%%%%%%%%%%%%%%%%%%%%%%%%%%%%%%%%%%%%%%%%%%%
% 1.4.6
%%%%%%%%%%%%%%%%%%%%%%%%%%%%%%%%%%%%%%%%%%%%%%%%
\subsection*{1.4.6 298.15 K における熱力学平衡定数}

平衡定数 \(K\) は
\[
 K = \exp\left(-\frac{\Delta_r G^\circ}{RT}\right)
\]
で与えられる．

\paragraph{(1) \(\mathrm{N_2(g)} + \mathrm{O_2(g)} \rightleftharpoons 2\mathrm{NO(g)}\)}

標準生成ギブズエネルギーは
\[
 \Delta_{\mathrm f}G^\circ(\mathrm{NO}) = +86.55\ \mathrm{kJ\,mol^{-1}},
\quad
 \Delta_{\mathrm f}G^\circ(\mathrm{N_2}) = \Delta_{\mathrm f}G^\circ(\mathrm{O_2}) = 0
\]
であるから，
\[
 \Delta_r G^\circ
 = 2\Delta_{\mathrm f}G^\circ(\mathrm{NO})
 = 2\times 86.55\ \mathrm{kJ\,mol^{-1}}
 = 173.10\ \mathrm{kJ\,mol^{-1}}.
\]
したがって
\[
 K
 = \exp\left(-\frac{173.10\times10^3}{(8.314)(298.15)}\right)
 \simeq 4.7\times10^{-31}.
\]
\[
 \boxed{K \simeq 4.7\times10^{-31}}
\]

\paragraph{(2) \(\dfrac12\mathrm{N_2(g)} + \dfrac32\mathrm{H_2(g)} \rightleftharpoons \mathrm{NH_3(g)}\)}

この反応は 1.4.4 で用いた反応と同じであり，
\(\Delta_r G^\circ = \Delta_{\mathrm f}G^\circ(\mathrm{NH_3}) = -16.45\ \mathrm{kJ\,mol^{-1}}\) である．
したがって
\[
 K
 = \exp\left(-\frac{-16.45\times10^3}{(8.314)(298.15)}\right)
 \simeq 7.6\times10^{2}.
\]
\[
 \boxed{K \simeq 7.6\times10^{2}}
\]


%%%%%%%%%%%%%%%%%%%%%%%%%%%%%%%%%%%%%%%%%%%%%%%%
% 1.5.1 ファントホッフの式と平衡定数の温度依存性
%%%%%%%%%%%%%%%%%%%%%%%%%%%%%%%%%%%%%%%%%%%%%%%%
\subsection*{1.5.1 ファントホッフの式の適用}

ファントホッフの式
\[
 \frac{\mathrm d\ln K}{\mathrm dT} = \frac{\Delta_r H^\circ}{RT^2}
\]
を温度について積分すると，
\[
 \ln\frac{K_2}{K_1}
 = -\frac{\Delta_r H^\circ}{R}\left(\frac{1}{T_2}-\frac{1}{T_1}\right)
\]
したがって
\[
 K_2
 = K_1 \exp\left[
        -\frac{\Delta_r H^\circ}{R}
         \left(\frac{1}{T_2}-\frac{1}{T_1}\right)
       \right]
\]
である．ここで \(T_1=298.15\ \mathrm{K}\) を基準温度とする．

\paragraph{(1) \(\dfrac12\mathrm{N_2} + \dfrac12\mathrm{O_2} \rightleftharpoons \mathrm{NO}\)}

この反応は NO の生成反応であるから，
\[
 \Delta_r H^\circ = \Delta_{\mathrm f}H^\circ(\mathrm{NO})
  = +90.25\ \mathrm{kJ\,mol^{-1}},
\quad
 \Delta_r G^\circ = \Delta_{\mathrm f}G^\circ(\mathrm{NO})
  = +86.55\ \mathrm{kJ\,mol^{-1}}.
\]
25\(^\circ\mathrm{C}\) における平衡定数は
\[
 K(298.15\ \mathrm{K})
 = \exp\left(-\frac{\Delta_r G^\circ}{RT_1}\right)
 = \exp\left(-\frac{86.55\times10^3}{(8.314)(298.15)}\right)
 \simeq 6.86\times10^{-16}.
\]

この \(K_1\) を用いて，各温度における平衡定数は
\[
 K(T_2)
 = K(298.15)\,
   \exp\left[
     -\frac{90.25\times10^3}{8.314}
      \left(\frac{1}{T_2}-\frac{1}{298.15}\right)
   \right]
\]
で与えられる．計算結果の一例を示すと，

\[
\begin{array}{c|c}
 T / \mathrm{K} & K \\
\hline
298.15 & 6.9\times10^{-16} \\
773.15\ (500^\circ\mathrm{C}) & 3.6\times10^{-6} \\
1773.15\ (1500^\circ\mathrm{C}) & 9.8\times10^{-3} \\
2273.15\ (2000^\circ\mathrm{C}) & 3.8\times10^{-2} \\
2773.15\ (2500^\circ\mathrm{C}) & 8.9\times10^{-2}
\end{array}
\]

となり，\(\Delta_r H^\circ>0\) の吸熱反応では，温度が上昇するほど
平衡定数が大きくなる（生成物側に有利になる）ことが分かる．

\paragraph{(2) \(\dfrac12\mathrm{N_2} + \dfrac32\mathrm{H_2} \rightleftharpoons \mathrm{NH_3}\)}

この反応では
\[
 \Delta_r H^\circ = -45.94\ \mathrm{kJ\,mol^{-1}},
\quad
 \Delta_r G^\circ = -16.45\ \mathrm{kJ\,mol^{-1}}.
\]
25\(^\circ\mathrm{C}\) における平衡定数は
\[
 K(298.15\ \mathrm{K})
 = \exp\left(-\frac{-16.45\times10^3}{(8.314)(298.15)}\right)
 \simeq 7.6\times10^{2}.
\]

同様に
\[
 K(T_2)
 = K(298.15)\,
   \exp\left[
     -\frac{-45.94\times10^3}{8.314}
      \left(\frac{1}{T_2}-\frac{1}{298.15}\right)
   \right]
\]
より，

\[
\begin{array}{c|c}
 T / \mathrm{K} & K \\
\hline
298.15 & 7.6\times10^{2} \\
773.15\ (500^\circ\mathrm{C}) & 8.7\times10^{-3} \\
1773.15\ (1500^\circ\mathrm{C}) & 1.5\times10^{-4} \\
2273.15\ (2000^\circ\mathrm{C}) & 7.7\times10^{-5} \\
2773.15\ (2500^\circ\mathrm{C}) & 5.0\times10^{-5}
\end{array}
\]

となる．このように，\(\Delta_r H^\circ<0\) の発熱反応では
温度が上がるほど平衡定数は小さくなり，生成物側ではなく
反応物側に有利になることが分かる．  

これらの結果を \(\ln K\) vs. \(1/T\) のグラフとしてプロットすると，
直線の傾きが \(-\Delta_r H^\circ/R\) であることが確認できる．


%%%%%%%%%%%%%%%%%%%%%%%%%%%%%%%%%%%%%%%%%%%%%%%%
% 1.5.2 平衡定数の温度依存性の違い
%%%%%%%%%%%%%%%%%%%%%%%%%%%%%%%%%%%%%%%%%%%%%%%%
\subsection*{1.5.2 平衡定数が温度で増減する理由}

ファントホッフの式
\[
 \frac{\mathrm d\ln K}{\mathrm dT}
 = \frac{\Delta_r H^\circ}{RT^2}
\]
から分かるように，平衡定数の温度依存性は
\(\Delta_r H^\circ\) の符号によって決まる．すなわち，

\begin{itemize}
  \item \(\Delta_r H^\circ > 0\)（吸熱反応）のとき：
    \[
      \frac{\mathrm d\ln K}{\mathrm dT} > 0
    \]
    となり，温度が上昇すると \(K\) は増加する
    （高温ほど生成物側に有利になる）．
  \item \(\Delta_r H^\circ < 0\)（発熱反応）のとき：
    \[
      \frac{\mathrm d\ln K}{\mathrm dT} < 0
    \]
    となり，温度が上昇すると \(K\) は減少する
    （高温ほど反応物側に有利になる）．
\end{itemize}

したがって，「温度上昇で平衡定数が大きくなる場合」と
「小さくなる場合」の違いは，\(\Delta_r H^\circ\) が
正か負かという \textbf{反応のエンタルピー変化の符号}に由来する．





\newpage



\subsubsection*{問題 2.1.1}
主量子数 $n=1$ のときの軌道半径（ボーア半径 $a_0$）を求める。
式： $r_n = \frac{\epsilon_0 h^2 n^2}{\pi m e^2}$
\[
a_0 = \frac{(8.854 \times 10^{-12} \text{ F/m}) \times (6.626 \times 10^{-34} \text{ J s})^2 \times 1^2}{\pi \times (9.109 \times 10^{-31} \text{ kg}) \times (1.602 \times 10^{-19} \text{ C})^2} \approx 5.29 \times 10^{-11} \text{ m} \quad (0.529\text{ \AA})
\]

\subsubsection*{問題 2.1.2}
水素原子におけるK殻 ($n=1$) とL殻 ($n=2$) のエネルギー差 $\Delta E$ を求める。
エネルギー準位の式： $E_n = -\frac{m e^4}{8 \epsilon_0^2 h^2 n^2}$
\[
\Delta E = E_2 - E_1 = \left( -\frac{m e^4}{8 \epsilon_0^2 h^2 \cdot 2^2} \right) - \left( -\frac{m e^4}{8 \epsilon_0^2 h^2 \cdot 1^2} \right) = \frac{m e^4}{8 \epsilon_0^2 h^2} \left( 1 - \frac{1}{4} \right) = \frac{3}{4} \left( \frac{m e^4}{8 \epsilon_0^2 h^2} \right)
\]
ここで、基底状態のエネルギー $E_1 \approx -13.6$ eV であるため、
\[
\Delta E = 13.6 \text{ eV} \times \left( 1 - \frac{1}{4} \right) = 13.6 \times 0.75 = 10.2 \text{ eV}
\]




\newpage


% 物理化学基礎 第8回　問題2.1.3, 2.2.1–2.2.3 解答例

\section*{問題 2.1.3 の解答}

角運動量量子数を $\ell$ とすると，その副殻に含まれる軌道の数は
\[
m_\ell=-\ell,-\ell+1,\dots,\ell-1,\ell
\]
より $2\ell+1$ 個であり，1 つの軌道にはスピンの向きの違いにより
2 個まで電子が入る（パウリの排他原理）ので，副殻に入る最大電子数は
\[
2(2\ell+1)
\]
となる．したがって

\begin{itemize}
  \item d 軌道：$\ell=2$ なので，$2(2\cdot2+1)=10$ 個
  \item f 軌道：$\ell=3$ なので，$2(2\cdot3+1)=14$ 個
  \item g 軌道：$\ell=4$ なので，$2(2\cdot4+1)=18$ 個
\end{itemize}

よって，求める最大電子数は
\[
\boxed{\text{d 軌道 }10\ \text{個},\quad \text{f 軌道 }14\ \text{個},\quad
       \text{g 軌道 }18\ \text{個}}
\]
である．

\section*{問題 2.2.1 の解答}

原子軌道への電子の入り方は，資料に示されたエネルギー準位の順序  
\[
1s\ll 2s<2p\ll 3s<3p\ll 4s<3d<4p\ll\cdots
\]
と，パウリの排他原理およびフントの規則に従う．また，d 軌道と s 軌道の
例外として Cu, Cr では 3d 軌道が半閉殻または閉殻になる配置が安定である
ことに注意する．:contentReference[oaicite:0]{index=0}

\begin{enumerate}
  \item Ar, Kr

  \begin{align*}
    \mathrm{Ar}(Z=18)&: 1s^22s^22p^63s^23p^6 = [\mathrm{Ne}]3s^23p^6,\\
    \mathrm{Kr}(Z=36)&: 1s^22s^22p^63s^23p^64s^23d^{10}4p^6
                      = [\mathrm{Ar}]3d^{10}4s^24p^6.
  \end{align*}

  \item Na, K

  \begin{align*}
    \mathrm{Na}(Z=11)&: 1s^22s^22p^63s^1 = [\mathrm{Ne}]3s^1,\\
    \mathrm{K}(Z=19)&: 1s^22s^22p^63s^23p^64s^1 = [\mathrm{Ar}]4s^1.
  \end{align*}

  \item Na$^+$, Mg$^{2+}$, Al$^{3+}$

  いずれも 10 個の電子をもつ等電子種であり，Ne と同じ配置になる．
  \begin{align*}
    \mathrm{Na}^+ &: [\mathrm{Ne}],\\
    \mathrm{Mg}^{2+} &: [\mathrm{Ne}],\\
    \mathrm{Al}^{3+} &: [\mathrm{Ne}].
  \end{align*}

  \item O$^{2-}$, F$^-$

  これらも 10 個の電子をもち，Ne と同じ配置である．
  \begin{align*}
    \mathrm{O}^{2-} &: [\mathrm{Ne}],\\
    \mathrm{F}^- &: [\mathrm{Ne}].
  \end{align*}

  \item Cu$^+$, Cu$^{2+}$

  中性原子 Cu の基底状態は，d 軌道が閉殻になる
  \[
    \mathrm{Cu}: [\mathrm{Ar}]3d^{10}4s^1
  \]
  である．イオン化するときは 4s 軌道の電子が先に失われるので，
  \begin{align*}
    \mathrm{Cu}^+ &: [\mathrm{Ar}]3d^{10},\\
    \mathrm{Cu}^{2+} &: [\mathrm{Ar}]3d^{9}.
  \end{align*}

  \item Cr

  Cr では，$3d^5$ の半閉殻が安定になるため，
  \[
    \mathrm{Cr}: [\mathrm{Ar}]3d^{5}4s^1
  \]
  が基底状態の電子配置である．
\end{enumerate}

\section*{問題 2.2.2 の解答}

典型元素（s ブロックおよび p ブロック元素）では，  
\[
\text{最外殻} = n s^{a} n p^{b} \quad (0 \le a \le 2,\ 0 \le b \le 6)
\]
のように，電子がすべて最外殻の s, p 軌道に入っており，
「特徴的な電子」は最外殻にある．

これに対して遷移元素（d ブロック・f ブロック元素）では，
最外殻の主量子数を $n$ とすると，電子は
\[
(n-1)d \ \text{あるいは}\ (n-2)f
\]
のような内殻の d, f 軌道にも入り，  
電子配置は概ね
\[
\cdots (n-1)d^{m}n s^{2}
\quad \text{あるいは}\quad
\cdots (n-2)f^{k}(n-1)d^{m}n s^{2}
\]
の形になる．すなわち，
\begin{itemize}
  \item 典型元素：\ 最外殻 ns, np 軌道が主役（ns$^a$np$^b$ 型）
  \item 遷移元素：\ 内殻の (n−1)d や (n−2)f 軌道が部分的に占有される
\end{itemize}
という違いがある．

\section*{問題 2.2.3 の解答}

d ブロックの遷移元素の基底状態は，多くの場合
\[
\cdots (n-1)d^{m}n s^2
\]
のように，最外殻 ns 軌道に 2 個の電子をもつ．  
資料にあるように d 軌道よりも $(n+1)s$ 軌道の方がわずかにエネルギー準位が低く，
イオン化の際にはエネルギーの高い ns 電子から取り去られやすい．:contentReference[oaicite:1]{index=1}

したがって，まず ns の 2 個の電子が失われて
\[
\cdots (n-1)d^{m}
\]
という配置になりやすく，多くの遷移元素が
\[
\boxed{2\ \text{価のカチオン}（M^{2+}）}
\]
をとりやすいのである．


\newpage

\section*{問題 2.3.1 の解答}

\subsection*{(1) 水素分子 \(\mathrm{H_2}\)}

水素原子の電子配置は
\[
1\mathrm{H}: 1s^1
\]
であり，価電子は \(1s\) 軌道に 1 個入っている。2 つの水素原子が近づくと，
結合軸（原子を結ぶ直線）に向かい合った 2 つの \(1s\) 軌道が同位相で重なり，
電子が 2 個入った結合性分子軌道をつくる。

原子価結合法の立場では，
\begin{itemize}
  \item 2 つの \(1s\) 軌道の重なりによる共有結合
  \item 結合軸に対して回転対称で，重なり領域は 1 箇所
\end{itemize}
であるから，この結合は
\[
\boxed{\text{\(1s\) 軌道どうしの \(\sigma\) 結合 1 本}}
\]
と表現される。:contentReference[oaicite:0]{index=0}

\subsection*{(2) フッ素分子 \(\mathrm{F_2}\)}

フッ素原子の電子配置は
\[
9\mathrm{F}: 1s^22s^22p^5
\]
であり，最外殻の \(2p\) 軌道に 5 個の電子が入っている。
\(2p_x,2p_y,2p_z\) のうち 2 つには電子が 2 個ずつ入り，
残りの 1 つの \(2p\) 軌道に不対電子が 1 個入る（図 2.3.4 参照）。

2 つのフッ素原子が結合軸を \(x\) 軸にして接近すると，
結合軸に沿った \(2p_x\) 軌道どうしが同位相で重なって共有結合をつくる。
この重なりは結合軸まわりに回転対称であり，
重なり領域は 1 箇所なので，
\[
\boxed{\text{\(2p\) 軌道どうしの \(\sigma\) 結合 1 本}}
\]
が形成される。残りの \(2p\) 軌道上の電子は
それぞれ孤立電子対として原子上に局在する。

\section*{問題 2.3.1 の解答}

\subsection*{(1) 水素分子 \(\mathrm{H_2}\)}

水素原子の電子配置は
\[
1\mathrm{H}: 1s^1
\]
であり，価電子は \(1s\) 軌道に 1 個入っている。2 つの水素原子が近づくと，
結合軸（原子を結ぶ直線）に向かい合った 2 つの \(1s\) 軌道が同位相で重なり，
電子が 2 個入った結合性分子軌道をつくる。

原子価結合法の立場では，
\begin{itemize}
  \item 2 つの \(1s\) 軌道の重なりによる共有結合
  \item 結合軸に対して回転対称で，重なり領域は 1 箇所
\end{itemize}
であるから，この結合は
\[
\boxed{\text{\(1s\) 軌道どうしの \(\sigma\) 結合 1 本}}
\]
と表現される。:contentReference[oaicite:0]{index=0}

\subsection*{(2) フッ素分子 \(\mathrm{F_2}\)}

フッ素原子の電子配置は
\[
9\mathrm{F}: 1s^22s^22p^5
\]
であり，最外殻の \(2p\) 軌道に 5 個の電子が入っている。
\(2p_x,2p_y,2p_z\) のうち 2 つには電子が 2 個ずつ入り，
残りの 1 つの \(2p\) 軌道に不対電子が 1 個入る（図 2.3.4 参照）。

2 つのフッ素原子が結合軸を \(x\) 軸にして接近すると，
結合軸に沿った \(2p_x\) 軌道どうしが同位相で重なって共有結合をつくる。
この重なりは結合軸まわりに回転対称であり，
重なり領域は 1 箇所なので，
\[
\boxed{\text{\(2p\) 軌道どうしの \(\sigma\) 結合 1 本}}
\]
が形成される。残りの \(2p\) 軌道上の電子は
それぞれ孤立電子対として原子上に局在する。


\section*{問題 2.3.3 の解答}

\subsection*{ダイヤモンド（ダイヤモンド結晶）}

ダイヤモンド中の炭素原子は 4 本の等価な共有結合をもつ。
炭素原子の \(2s\) 軌道 1 個と \(2p\) 軌道 3 個を平均化し，
4 つの \(\mathrm{sp^3}\) 混成軌道をつくる（図 2.3.8）。:contentReference[oaicite:2]{index=2}
各 \(\mathrm{sp^3}\) 混成軌道が周囲の 4 つの炭素原子と
\(\sigma\) 結合をつくることにより，
正四面体的な三次元網目構造が形成される。

したがって，
\[
\boxed{\text{ダイヤモンド中の C 原子は \(\mathrm{sp^3}\) 混成軌道}}
\]
をとると考えればよい。

\subsection*{グラファイト（黒鉛）}

グラファイトでは，炭素原子は
同一平面内で 3 本の結合をなし，六角形網目をつくる。
ここでは \(2s\) 軌道 1 個と \(2p\) 軌道 2 個を平均化して
3 つの \(\mathrm{sp^2}\) 混成軌道をつくり（図 2.3.9），
それぞれが同一平面内で 3 本の \(\sigma\) 結合を形成する。残った
1 つの \(2p_z\) 軌道（平面に垂直な \(p\) 軌道）が隣接する炭素原子の
\(2p_z\) と重なり，平面の上下に広がる \(\pi\) 結合系をつくる（図 2.3.10）。:contentReference[oaicite:3]{index=3}

したがって，
\[
\boxed{\text{グラファイト中の C 原子は \(\mathrm{sp^2}\) 混成軌道}}
\]
と考えればよい。

\section*{問題 2.3.4 の解答}

共鳴構造は，ベンゼンの例（図 2.3.11, 2.3.12）に見られるように，
\(\pi\) 電子が 2 原子間に局在せず，多数の原子にまたがって
非局在化しているときに導入される概念である。:contentReference[oaicite:4]{index=4}

ダイヤモンドでは，各 C 原子は \(\mathrm{sp^3}\) 混成軌道からの
\(\sigma\) 結合のみをもち，電子は結合軸付近に局在している。
\(\pi\) 電子系が存在しないので，共鳴構造を考える余地はほとんどない。

一方，グラファイトでは，
\(\mathrm{sp^2}\) 混成軌道による平面内の \(\sigma\) 結合に加えて，
平面に垂直な \(2p_z\) 軌道どうしの重なりによる
広く非局在化した \(\pi\) 電子系が存在する。
この \(\pi\) 電子は平面全体にわたって動き回ることができ，
ベンゼンと同様に多くの極限構造の共鳴として表現できる。

したがって，
\[
\boxed{\text{共鳴構造が考えられるのはグラファイト（黒鉛）であり，
ダイヤモンドではない。}}
\]

\newpage


%%%%%%%%%%%%%%%%%%%%%%%%%%%%%%%%%%%%%%%%%%%%%%%%%%%%%%%%%%%%
\section*{問題 2.4.1 He$_2$ 分子が安定に存在しないことを分子軌道法によって説明せよ。}

\subsection*{解答}

He 原子 $1$ 個の電子配置は
\[
\mathrm{He}: 1s^2
\]
であるから，He$_2$ 分子を考えると，$1s$ 原子軌道 (AO) が 2 個
結合して，図 2.4.2 のように結合性分子軌道 $\sigma_{1s}$ と
反結合性分子軌道 $\sigma_{1s}^*$ の 2 つの分子軌道 (MO) が生じる．:contentReference[oaicite:0]{index=0}

He$_2$ 全体では電子は $4$ 個であり，パウリの排他原理に従って
エネルギーの低い順に MO を占有させると
\[
\sigma_{1s}^2\,\sigma_{1s}^{*\,2}
\]
となる．

このとき結合性 MO に入る電子数を $N_\mathrm{b}$,
反結合性 MO に入る電子数を $N_\mathrm{a}$ とすると，
結合次数 $B$ は講義ノートにあるように
\[
B=\frac{N_\mathrm{b}-N_\mathrm{a}}{2}
\]
で与えられる．:contentReference[oaicite:1]{index=1}

He$_2$ では
\[
N_\mathrm{b}=2,\qquad N_\mathrm{a}=2
\]
であるから
\[
B=\frac{2-2}{2}=0
\]
となる．

結合次数が $0$ であることは，
He$_2$ が 2 個の He 原子に比べて安定化されないことを意味し，
そのため He$_2$ 分子は安定に存在しない。



%%%%%%%%%%%%%%%%%%%%%%%%%%%%%%%%%%%%%%%%%%%%%%%%%%%%%%%%%%%%
\section*{問題 2.4.2 He$_2^+$ が存在できることを説明せよ。}

\subsection*{解答}

He$_2^+$ では He$_2$ から電子を 1 個取り去った状態なので，
全電子数は $3$ 個である．よって MO への電子配置は
\[
\sigma_{1s}^2\,\sigma_{1s}^{*\,1}
\]
となる．

このとき
\[
N_\mathrm{b}=2,\qquad N_\mathrm{a}=1
\]
なので結合次数は
\[
B=\frac{N_\mathrm{b}-N_\mathrm{a}}{2}
 =\frac{2-1}{2}=\frac{1}{2}
\]
となる．

結合次数が $1/2>0$ であることから，
He$_2^+$ は 2 個の He 原子 (He + He$^+$) よりも
わずかに安定化される。
したがって He$_2^+$ は弱い結合ではあるが
分子として存在することができる。



%%%%%%%%%%%%%%%%%%%%%%%%%%%%%%%%%%%%%%%%%%%%%%%%%%%%%%%%%%%%
\section*{問題 2.4.3 N$_2$ 分子の分子軌道は図 2.4.5 のようになる。結合次数を求めよ。}

\subsection*{解答}

N 原子 1 個の電子配置は
\[
\mathrm{N}: 1s^2 2s^2 2p^3
\]
であり，価電子は $2s^2 2p^3$ の 5 個である。
したがって N$_2$ 分子全体の価電子数は $10$ 個となる。

図 2.4.5 のエネルギー準位図では，
2s 軌道から生じる $\sigma_1$，$\sigma_2^*$ と，
2p 軌道から生じる $\pi_{2p}$（縮退 2 本），$\sigma_3$，
$\pi_{2p}^*$（縮退 2 本），$\sigma_4^*$ の MO が示されている。:contentReference[oaicite:2]{index=2}

エネルギーの低い順に 10 個の電子を MO に入れると
\[
\sigma_1^2\,\sigma_2^{*\,2}\,
\pi_{2p}^4\,\sigma_3^2
\]
となり，反結合性 MO である $\pi_{2p}^*$，$\sigma_4^*$ は空軌道である。

結合性 MO に入った電子数は
\[
N_\mathrm{b} = 
2 (\sigma_1) +
4 (\pi_{2p}) +
2 (\sigma_3)
=8,
\]
反結合性 MO に入った電子数は
\[
N_\mathrm{a}=2 \quad (\sigma_2^*)
\]
であるから，結合次数は
\[
B=\frac{N_\mathrm{b}-N_\mathrm{a}}{2}
  =\frac{8-2}{2}=3
\]
となる。

ゆえに N$_2$ 分子の結合次数は $3$ であり，
共有結合は三重結合である。



\newpage



\section*{問題 2.5.1 解答}

共通の方針：sp$^2$ 炭素 1 個あたり $\pi$AO（$p_z$ 軌道）を 1 つ，
$\pi$ 電子を 1 個と数える。  
ヒュッケル法では
\[
x = \frac{\alpha-\varepsilon}{\beta}
\]
とおき，隣接する炭素同士の行列要素を $1$，それ以外を $0$ とした
永年行列式を立てる。

\subsection*{(1) 1,3-ブタジエン}

構造： \(\mathrm{CH_2\!=CH\!-\!CH\!=CH_2}\) （直鎖共役二重結合）

\begin{itemize}
  \item AO 数：共役に関与する C が 4 個 $\Rightarrow$ 4 AO
  \item $\pi$ 電子数：二重結合 2 本 $\Rightarrow 4$ 電子
\end{itemize}

C 原子を左から 1,2,3,4 と番号付けすると，永年行列式は
\[
\left|
\begin{array}{cccc}
x & 1 & 0 & 0 \\
1 & x & 1 & 0 \\
0 & 1 & x & 1 \\
0 & 0 & 1 & x
\end{array}
\right| = 0.
\]

\subsection*{(2) 1,3-シクロブタジエン}

構造：正方形環 \(\mathrm{C_4H_4}\) に 1,3-位で二重結合
（4 つの C が環状にすべて隣接）

\begin{itemize}
  \item AO 数：共役 C が 4 個 $\Rightarrow$ 4 AO
  \item $\pi$ 電子数：二重結合 2 本 $\Rightarrow 4$ 電子
\end{itemize}

C 原子を環に沿って 1,2,3,4 と番号付けると，
隣接結合は 1–2, 2–3, 3–4, 4–1 であるから永年行列式は
\[
\left|
\begin{array}{cccc}
x & 1 & 0 & 1 \\
1 & x & 1 & 0 \\
0 & 1 & x & 1 \\
1 & 0 & 1 & x
\end{array}
\right| = 0.
\]

\subsection*{(3) プロペニルラジカル（アリルラジカル）}

構造： \(\mathrm{CH_2\!=CH\!-\!CH_2^{\boldsymbol{\cdot}}}\)  
（3 個の sp$^2$ C が直鎖状に共役したラジカル）

\begin{itemize}
  \item AO 数：共役 C が 3 個 $\Rightarrow$ 3 AO
  \item $\pi$ 電子数：$\pi$ 結合 1 本で 2 電子 + ラジカルの 1 電子
        $\Rightarrow 3$ 電子
\end{itemize}

C 原子を左から 1,2,3 とすると，隣接結合は 1–2, 2–3 なので
永年行列式は
\[
\left|
\begin{array}{ccc}
x & 1 & 0 \\
1 & x & 1 \\
0 & 1 & x
\end{array}
\right| = 0.
\]




\newpage




%%%%%%%%%%%%%%%%%%%%%%%%%%%%%%%%%%%%%%%%%%%%%%%%%%
% 物理化学基礎 2025 第12回 ヒュッケル法 演習解答
%%%%%%%%%%%%%%%%%%%%%%%%%%%%%%%%%%%%%%%%%%%%%%%%%%

\section*{問題 2.5.2 プロペニルカチオン（アリルカチオン）}

分子式 CH$_2$=CH-CH$_2^+$ を考える。$\pi$ 電子に関与するのは
sp$^2$ 混成した C 原子 3 個の $p$ 軌道であり，$\pi$ 電子数は
\[
3\times 1 - 1 = 2
\]
個である（カチオンなので 1 個減る）。

\subsection*{(1) 永年行列式と $x$ の値}

原子 1,2,3 を C-C-C の順にとると，共役は 1-2 および 2-3 のみである。
したがってヒュッケル行列は
\[
\begin{vmatrix}
x & 1 & 0 \\
1 & x & 1 \\
0 & 1 & x
\end{vmatrix}=0
\]
である。行列式を展開すると
\[
x(x^2-1) - 1\cdot x = x^3-2x = 0
\]
より
\[
x(x^2-2)=0 \quad\Rightarrow\quad
x = 0,\ \pm\sqrt{2}.
\]

\subsection*{(2) MO エネルギーと全電子エネルギー}

ヒュッケル法では MO エネルギーは
\[
\varepsilon = \alpha - x\beta
\]
で与えられる（$\beta<0$）。したがって 3 つの MO のエネルギーは
\begin{align*}
x=-\sqrt{2}&:\quad \varepsilon_1 = \alpha + \sqrt{2}\beta,\\
x=0       &:\quad \varepsilon_2 = \alpha,\\
x=+\sqrt{2}&:\quad \varepsilon_3 = \alpha - \sqrt{2}\beta.
\end{align*}
$\beta<0$ なので $x$ がもっとも負である $x=-\sqrt{2}$ の軌道
$\varphi_1$ が最も安定である。

アリルカチオンの $\pi$ 電子は 2 個なので，$\varphi_1$ に 2 電子を
入れればよく，全電子エネルギーは
\[
E_{\mathrm{tot}} = 2\varepsilon_1
                  = 2(\alpha + \sqrt{2}\beta)
\]
となる。数値が必要なら，資料と同じ
$\alpha=-7.06\ \mathrm{eV},\ \beta=-2.49\ \mathrm{eV}$ を用いて
\[
\varepsilon_1 \approx -10.6\ \mathrm{eV},\qquad
E_{\mathrm{tot}} \approx -21.2\ \mathrm{eV}
\]
となる。

\subsection*{(3) 係数 $C_i$ と電子密度}

$x=-\sqrt{2}$ に対応する永年方程式は
\[
\begin{cases}
xC_1 + C_2      = 0,\\[2pt]
C_1 + xC_2 + C_3 = 0,\\[2pt]
C_2 + xC_3      = 0,
\end{cases}
\qquad (x=-\sqrt{2})
\]
である。第1式から
\[
C_2 = \sqrt{2}\,C_1,
\]
第3式から
\[
\sqrt{2}\,C_1 + (-\sqrt{2})C_3 = 0 \quad\Rightarrow\quad C_3 = C_1
\]
となり，第2式も自動的に満たされる。したがって固有ベクトルは
\[
(C_1,C_2,C_3) \propto (1,\sqrt{2},1)
\]
である。規格化条件
\[
|C_1|^2+|C_2|^2+|C_3|^2 = 1
\]
より
\[
C_1=C_3=\frac12,\qquad
C_2=\frac{\sqrt{2}}{2}
\]
となる（全体にマイナスを掛けたものも同等）。

MO $\varphi_1$ に電子 2 個が入っているので，原子 $i$ 上の
$\pi$ 電子密度 $\rho_i$ は
\[
\rho_i = 2|C_i|^2
\]
で与えられる。よって
\[
\rho_1 = \rho_3 = 2\left(\frac12\right)^2 = 0.5,\qquad
\rho_2 = 2\left(\frac{\sqrt{2}}2\right)^2 = 1.0.
\]

\subsection*{(4) 各原子の電荷}

1 個の C 原子について $\pi$ 電子 1 個だけを考えているので，
C の正電荷は +1 とみなす（資料中のシクロプロペニルカチオンと同じ扱い）。
したがって各原子の電荷 $q_i$ は
\[
q_i = +1 - \rho_i
\]
であり，
\[
q_1 = q_3 = +1 - 0.5 = +0.5,\qquad
q_2 = +1 - 1.0 = 0
\]
となる。

すなわち，正電荷（カチオン性）は端の 2 つの C 原子に
それぞれ +0.5 ずつ分布し，中央の C 原子はほぼ中性であると
予想される。

%%%%%%%%%%%%%%%%%%%%%%%%%%%%%%%%%%%%%%%%%%%%%%%%%%
\section*{問題 2.5.3 プロトン化水素 H$_3^+$ の構造}

H$_3^+$ の 3 個の 1s 軌道に対しても，同様にヒュッケル型の取り扱いを行う。
価電子数は 2 個であり，C$_3$H$_3^+$ の場合と同じく
「3 原子 2 電子系」である。

\subsection*{(1) 直線構造の場合}

直線構造 H-H-H を考えると，結合は 1-2 および 2-3 のみであり，
アリルカチオンと同じトポロジーになる。したがって
永年行列式は
\[
\begin{vmatrix}
x & 1 & 0 \\
1 & x & 1 \\
0 & 1 & x
\end{vmatrix}=0
\]
であり，固有値は
\[
x = 0,\ \pm\sqrt{2}.
\]
エネルギーは
\[
\varepsilon = \alpha - x\beta
\]
より
\[
\varepsilon_1 = \alpha + \sqrt{2}\beta,\quad
\varepsilon_2 = \alpha,\quad
\varepsilon_3 = \alpha - \sqrt{2}\beta.
\]
2 電子なので最下位 MO のみに 2 電子が入り，
全電子エネルギーは
\[
E_{\mathrm{lin}} = 2(\alpha + \sqrt{2}\beta).
\]

\subsection*{(2) 正三角形構造の場合}

3 個の H が正三角形に配置されると，3 原子はすべて互いに隣接しており，
シクロプロペニルカチオンと同じトポロジーになる。よって永年行列式は
\[
\begin{vmatrix}
x & 1 & 1 \\
1 & x & 1 \\
1 & 1 & x
\end{vmatrix}=0
\]
であり，
\[
x = -2,\ +1,\ +1
\]
を得る。エネルギーは
\[
\varepsilon_1 = \alpha - (-2)\beta = \alpha + 2\beta,\quad
\varepsilon_2 = \varepsilon_3 = \alpha - \beta
\]
であり，2 電子は最も低い軌道 $\varepsilon_1$ に入るので
\[
E_{\triangle} = 2(\alpha + 2\beta)
\]
となる。

\subsection*{(3) 安定な構造の比較}

$\beta<0$ であることに注意すると
\[
\alpha + 2\beta\ <\ \alpha + \sqrt{2}\beta
\]
が成り立つ（$2 > \sqrt{2}$ なので，負の数 $\beta$ を掛けると順序が反転する）。
したがって
\[
E_{\triangle} = 2(\alpha + 2\beta)
             < 2(\alpha + \sqrt{2}\beta) = E_{\mathrm{lin}}
\]
となり，正三角形構造の方が直線構造よりも
全電子エネルギーが低く，より安定である。

よって，ヒュッケル法に基づけば
\[
\mathrm{H_3^+}
\]
は直線構造ではなく \textbf{正三角形構造} をとると予想される。

\newpage

%%%%%%%%%%%%%%%%%%%%%%%%%%%%%%%%%%%%%%%%%%%%%%%%%%%%%%%%%%%%
\section*{問題 2.5.4 ベンゾ[a]アントラセンのフロンティア軌道}

ベンゾ[a]アントラセン C$_{18}$H$_{12}$ は，$\pi$ 電子を 18 個もつ。
前頁の表では分子軌道 $\phi_i$（$i=1,\dots ,18$）のエネルギー $\varepsilon_i$ が
低い順に
\[
\varepsilon_1<\varepsilon_2<\cdots<\varepsilon_{18}
\]
として与えられている。

\subsection*{(1) HOMO および LUMO}

1 つの分子軌道にはスピンの違う 2 個までの電子が入るので，
18 個の $\pi$ 電子は
\[
\phi_1,\phi_2,\dots,\phi_9
\]
の 9 個の軌道を 2 電子ずつ満たす。したがって
\begin{itemize}
  \item 最高被占軌道（HOMO）は $\phi_9$，
  \item 最低空軌道（LUMO）は $\phi_{10}$
\end{itemize}
である。

表よりそれぞれのエネルギーは
\[
\varepsilon_9 = -8.186\ {\rm eV},\qquad
\varepsilon_{10} = -5.934\ {\rm eV}
\]
である。

\subsection*{(2) 全 $\pi$ 電子エネルギー}

全 $\pi$ 電子エネルギー $E_\pi$ は，被占軌道のエネルギーの 2 倍和
\[
E_\pi = 2\sum_{i=1}^{9} \varepsilon_i
\]
で与えられる。表の数値を用いると
\begin{align*}
\sum_{i=1}^{9}\varepsilon_i
&=\varepsilon_1+\varepsilon_2+\cdots+\varepsilon_9 \\
&= (-13.247)+(-12.477)+(-11.429)+(-10.745)+(-10.354) \\
&\hspace{1em} +(-9.962)+(-9.550)+(-8.840)+(-8.186) \\
&= -94.79\ {\rm eV}
\end{align*}
したがって
\[
E_\pi = 2\times (-94.79)\ {\rm eV}
      = -1.90\times 10^2\ {\rm eV}\quad(\simeq -189.6\ {\rm eV})
\]
となる。

\subsection*{(3) 求電子反応に対して最も反応性の高い炭素と，
そのときの HOMO の電子密度}

求電子反応では，フロンティア軌道理論より，
HOMO 上の電子密度が大きい炭素ほど反応性が高いと考えられる。

HOMO は $\phi_9$ であり，
原子 $j$ 上の $\phi_9$ の\emph{電子密度} $\rho_j^{\rm(HOMO)}$ は
（$\phi_9$ に 2 電子入っているので）
\[
\rho_j^{\rm(HOMO)} = 2\,C_{9j}^2
\]
で与えられる。ここで $C_{9j}$ は $\phi_9$ の原子軌道展開係数である。

前頁の表から $\phi_9$ の係数を読むと
\[
(C_{9,1},C_{9,2},\dots,C_{9,18}) =
(0.003,\,0.204,\,-0.078,\,-0.393,\,-0.100,\,0.301,\,0.236,\,-0.194,
\]
\[
\hspace{7.3em}
-0.324,\,0.047,\,0.445,\,0.154,\,-0.298,\,-0.288,\,0.167,\,0.160,\,-0.095,\,-0.203)
\]
であるから，それぞれの電子密度は
\[
\rho_j^{\rm(HOMO)} = 2C_{9j}^2
\]
を計算すればよい。例えば $j=11$ では
\[
\rho_{11}^{\rm(HOMO)}
=2\times(0.445)^2
\approx 0.40
\]
となる。同様に他の炭素について計算すると，
$\rho_j^{\rm(HOMO)}$ の最大値は炭素 11 番で得られる
（次いで 4 番, 9 番, 6 番 \dots の順に小さくなる）。

したがって
\begin{itemize}
  \item 求電子反応に対して最も反応性の高い炭素は 11 番炭素，
  \item そのときの HOMO における電子密度は
        \[
        \rho_{11}^{\rm(HOMO)} \simeq 0.40
        \]
\end{itemize}
と結論される。








\end{document}